\documentclass[11pt]{article}
\usepackage[margin=1in]{geometry}
\usepackage[T1]{fontenc}
\usepackage[utf8]{inputenc}
\usepackage{lmodern,microtype,amsmath,amssymb,amsthm,mathtools,needspace}
\usepackage[colorlinks=true,linkcolor=blue,citecolor=blue,urlcolor=blue]{hyperref}
\hypersetup{pdftitle={Sparse parity under globally conditioned neural priors},pdfauthor={ChatGPT-6 Astra; Kaarel Hanni; Dmitry Vaintrob}}
\setlength{\emergencystretch}{2em}
\newtheorem{theorem}{Theorem}[section]
\newtheorem{lemma}[theorem]{Lemma}
\newtheorem{proposition}[theorem]{Proposition}
\newtheorem{corollary}[theorem]{Corollary}
\theoremstyle{remark}\newtheorem{remark}[theorem]{Remark}
\DeclareMathOperator{\Var}{Var}
\DeclareMathOperator{\sign}{sign}
\DeclareMathOperator{\dist}{dist}
\DeclareMathOperator{\KL}{KL}
\newcommand{\E}{\mathbb E}
\newcommand{\Pp}{\mathbb P}
\title{Sparse parity under globally conditioned neural priors}
\author{ChatGPT-6 Astra \and Kaarel H\"anni \and Dmitry Vaintrob\thanks{Authors are listed in alphabetical order.}}
\date{October 15, 2026}
\begin{document}
\maketitle
\begin{abstract}
We study posterior confidence for a parity of $k$ coordinates in a
$d$-dimensional Boolean input under a specified Bayesian neural-network
prior. The network has fixed depth, Gaussian weights, hyperbolic-tangent
hidden activations, and a linear readout. Its prior is conditioned to have
a fixed absolute output margin on every input, followed by exact training
labels. For every fixed negative width-scaling exponent $\gamma<0$, we
prove that each sample size
$n\le\exp\{c k\log(ed/k)\}$ leaves a universal positive fraction of
inputs without correct-label posterior probability greater than $2/3$,
with probability at least $3/4$, in the width-first limit. This is a
$d^{\Omega(k)}$ lower bound when $k\le d^{1-\eta}$ for fixed $\eta>0$.
The proof uses Gaussian sign conditioning, sharp Fourier spectra,
near-orthogonal covariance bounds, and a uniform truncated negative-part
estimate. The distinction between scaling regimes is substantive:
for every $0<\gamma<1$, the same prior learns a one-coordinate parity
with $O(\log d)$ samples. We prove this by identifying its exact minimum
energy sign functions. These results do not settle the corresponding
growing-$k$ lower bound at $\gamma=0$ or $0<\gamma<1$.
\end{abstract}

\section{Model and results}
Let $\mathcal X=\{-1,1\}^d$ with uniform probability measure.
Fix $U\subseteq[d]$, $|U|=k$, and set
$f(x)=\chi_U(x)=\prod_{j\in U}x_j$.
All hidden layers have width $N$, the number of weight matrices
$L\ge2$ is fixed, and
\[
 h^0(x)=(1,x),\qquad
 h^s(x)=\tanh(W^sh^{s-1}(x))\quad(1\le s<L),
 \qquad z_N(x)=W^Lh^{L-1}(x).
\]
There are no additional biases. All weights are independent centered
Gaussians, with variances
\begin{equation}\label{eq:weights}
 \Var W^1_{ij}=N^{1-\gamma}/d,\qquad
 \Var W^s_{ij}=N^{-\gamma}\quad(1<s<L),\qquad
 \Var W^L_i=N^{-\gamma-1}.
\end{equation}
Fix $a>0$; in particular, one may take $a=1/2$.
Condition the prior on $|z_N(x)|>a$ for every $x\in\mathcal X$.
For independent uniform training inputs $S=(X_1,\ldots,X_n)$,
condition subsequently on $f(X_i)z_N(X_i)>0$ for every $i$.
Equivalently, condition on the intersection of these events.
Define
\begin{align*}
 q_{N,S}(x)&=\mathbb P\{f(x)z_N(x)>0\mid
 |z_N(y)|>a\ \forall y;\ f(X_i)z_N(X_i)>0\ \forall i\},\\
 F_N(S)&=\mathbb P_x\{q_{N,S}(x)\le2/3\}.
\end{align*}
Thus a successful prediction requires strictly more than $2/3$
posterior probability of the correct label. A fresh input is always
independent and uniform. The learner is not told the parity support.

\Needspace{12\baselineskip}
\begin{theorem}[Sparse-parity sample lower bound]\label{thm:main}
For each fixed $L\ge2$ and $a>0$, there are $c>0$ and $d_0$
such that, for every $d\ge d_0$, every $1\le k\le d$, every
support $U$ of size $k$, every integer
\[
 0\le n\le\exp\!\left(c k\log\frac{ed}{k}\right),
\]
and every fixed $\gamma<0$,
\[
 \mathbb P_S\left\{\liminf_{N\to\infty}F_N(S)
          \ge\varepsilon\right\}\ge\frac34,
 \qquad
 \varepsilon=\frac{2\pi(1/e-7/20)^2}{256}>0.
\]
The conditioning events have positive probability for all sufficiently
large widths at each fixed dimension and sample.
\end{theorem}

The numerical value of the displayed failure fraction is approximately
$7.85\times10^{-6}$. It is uniform in $d,k,L,a$, although the sample
exponent and dimension threshold are not quantitatively optimized.
In particular, the theorem rules out vanishing confidence failure below
the stated sample scale; it does not assert a failure fraction as large
as $1/10$.
For any fixed $\eta\in(0,1)$ and $k\le d^{1-\eta}$, the bound
includes every $n\le d^{c\eta k}$. When $k=d$ it becomes
$n\le e^{cd}$. The latter dependence is necessary for a uniform
statement: the input space has only $2^d$ points, and observing almost
all of them suffices for this confidence criterion under noiseless
conditioning.

The width limit is taken first, at fixed $d,k,L,a,\gamma,S$.
No simultaneous dimension--width limit or quantified minimum width is
asserted. The probability bound holds for each specified sample size
with uniform constants, rather than on a common event for all prefixes
of an infinite sample sequence. It uses no monotonicity assumption on
fixed-target posterior risk.

\begin{theorem}[A separation between scaling regimes]
\label{thm:separation}
Fix $d\ge1$, $L\ge2$, $a>0$, and $0<\delta<1$.
For $0<\gamma<1$, a target $f(x)=x_1$ and
\[
 n\ge\left\lceil\log_2\frac{2d}{\delta}\right\rceil
\]
satisfy $q_{N,S}(x)\to1$ for every input, with probability at least
$1-\delta$ over training samples. In contrast, for $\gamma<0$,
every target and every sample of size $n$ satisfy
\[
 \liminf_{N\to\infty}F_N(S)
 \ge1-\left(2+\frac{400}{9\pi}\right)\frac{n}{d+1}.
\]
\end{theorem}

Thus a $d^{c k}$ lower bound cannot hold uniformly over all
$\gamma<1$ and all $k\ge1$.
Theorem~\ref{thm:separation} does not resolve the lower-bound question
restricted to growing $k$ for $0\le\gamma<1$.

\subsection{Proof overview}
Put $\Psi(t)=2\arcsin(t)/\pi$ and $\beta=2^{-(L-1)}$.
The normalized network outputs converge to a Gaussian vector with
covariance
\[
 K(x,y)=\Psi^{\circ(L-1)}\!\left(\frac{1+x\cdot y}{d+1}\right).
\]
We regard $K$ as a $2^d$-by-$2^d$ covariance matrix. Its normalized
operator is $2^{-d}K$, and the kernel norm of a field $h$ is
$\|h\|_K^2=h^{\mathsf T}K^{-1}h$. Population norms, in contrast,
use the uniform probability measure on the cube. Keeping these two
normalizations separate is essential in the spectral and sample bounds.
For $\gamma<0$, the fixed global margin vanishes after normalization,
leaving Gaussian conditioning on the training signs. A homogeneous-cone
argument bounds the energy of its actual posterior mean by $n$.
Log-concavity and Gaussian regression convert posterior confidence
into a quantitative mean margin wherever residual variance remains.
These are implications about the actual posterior, not a replacement
of its predictions by the sign of its mean.

We split that mean into a low-degree part, one further pair of Fourier
levels, and a remaining tail. Uniform sampling controls a truncated
negative part of the first two pieces. Orthogonality to the target
parity forces their positive and negative integrals to agree whenever
their degree is below $k$. Sharp covariance estimates control the tail
and prevent the negative integral from being hidden on a tiny set.
A near-orthogonal entry bound is what permits cutoffs much smaller than
$\log d$. This yields the ambient-dimension exponent uniformly in $k$.

For $0<\gamma<1$, global-margin conditioning instead selects compatible
sign functions of minimum Gaussian energy. The top eigenspace of the
kernel consists exactly of affine functions, and the only globally
constant-magnitude functions in it are constants and signed coordinates.
This identifies the ground states and proves the logarithmic upper bound.
All required width-limit, spectral, sampling, and symmetry arguments are
included below.

\section{The kernel rate of the actual network}
Let $M=2^d$, and identify output functions with vectors in $\mathbb R^M$.
Define
\begin{equation}\label{eq:kernel}
 H(t)=\frac2\pi\arcsin t,\qquad
 K(x,y)=H^{\circ(L-1)}\!\left(\frac{1+x\cdot y}{d+1}\right).
\end{equation}
The positive definiteness of $K$ is proved in the next section.
For each Boolean function $g$, put
\begin{equation}\label{eq:energy}
 E_K(g)=\min_{g(x)z(x)\ge a\ \forall x}\frac12 z^TK^{-1}z.
\end{equation}
The minimum exists since its feasible set is closed and nonempty and
the objective is coercive. Replacing the closed inequalities by strict
ones gives the same infimum, by scaling a feasible vector by $1+\epsilon$.

\begin{proposition}[Posterior concentration on minimum-energy signs]
\label{prop:rate}
For the actual network~\eqref{eq:weights}, at every fixed
$0<\gamma<1$,
\begin{equation}\label{eq:rate}
 \lim_{N\to\infty}N^{-\gamma}\log
 \mathbb P\{g(x)z_N(x)>a\ \forall x\}=-E_K(g).
\end{equation}
Consequently, after any fixed noiseless training sample, posterior
mass tends to one on the compatible functions minimizing $E_K$.
No assertion about the relative weights of tied minimizers is needed.
\end{proposition}

\begin{proof}
For hidden layer $s$, let
$Q_N^s(x,y)=N^{-1}\sum_i h_i^s(x)h_i^s(y)$.
We first establish, for each fixed $\epsilon>0$,
\begin{equation}\label{eq:freeze}
 \mathbb P\{\|Q_N^{L-1}-K\|>\epsilon\}
       \le C_\epsilon e^{-c_\epsilon N}
\end{equation}
for all sufficiently large $N$. Any matrix norm can be used, since
$M$ is fixed. Conditional on the preceding layer, the rows of the
next hidden layer are independent bounded vectors. Hoeffding's
inequality controls all their empirical second moments. Their
preactivations have covariance $N^{1-\gamma}Q_N^{s-1}$; at the first
layer the covariance is
$N^{1-\gamma}(1+x\cdot y)/d$. Because $N^{1-\gamma}\to\infty$,
the conditional second-moment map converges to the Gaussian-sign map
\[
 Q\longmapsto \left[H\!\left(
 \frac{Q(x,y)}{\sqrt{Q(x,x)Q(y,y)}}\right)\right]_{x,y}.
\]
This follows by dominated convergence for a Gaussian vector, since
each scalar variance stays bounded away from zero in a neighborhood
of the limiting covariance. The convergence is uniform in such a
neighborhood: a compactness subsequence and a continuous square root
coupling reduce the assertion to the same dominated-convergence
argument. The limiting map is continuous there. Inducting through
the fixed number of layers, with smaller fixed tolerances at earlier
layers, proves~\eqref{eq:freeze}. The first-layer limit normalizes
the diagonal $(d+1)/d$, giving the kernel in~\eqref{eq:kernel}.

Conditionally on the hidden layers, the readout is exactly Gaussian:
\begin{equation}\label{eq:readout}
 Z_N\mid Q_N^{L-1}\sim\mathcal N(0,N^{-\gamma}Q_N^{L-1}).
\end{equation}
Couple it to $Y_N=N^{-\gamma/2}K^{1/2}\xi$, where
$\xi\sim\mathcal N(0,I_M)$ is independent of the hidden layers.
On a fixed sufficiently small neighborhood of the positive-definite
matrix $K$, the square-root map is Lipschitz. Hence, for $\delta>0$,
\[
 \mathbb P\{\|Z_N-Y_N\|>\delta\}
 \le C_\epsilon e^{-c_\epsilon N}
 +\mathbb P\{\|\xi\|>\delta N^{\gamma/2}/(C\epsilon)\}.
\]
First divide logarithms by $N^\gamma$ and let $N\to\infty$;
then let $\epsilon\downarrow0$. The resulting upper limit is
$-\infty$, since $\gamma<1$. The coupled vectors are therefore
exponentially equivalent at speed $N^\gamma$.

The fixed-dimensional Gaussian vectors $Y_N$ have large-deviation
rate $I(z)=z^TK^{-1}z/2$, directly from their densities. Exponential
equivalence transfers the same bounds to $Z_N$. Applying the upper
bound to a closed margin orthant and the lower bound to its interior
gives~\eqref{eq:rate}, since their infima agree. Finally, there are
only finitely many Boolean functions for fixed $d$. Dividing each
compatible orthant probability by their sum proves the posterior
concentration claim.
\end{proof}

This is a rate statement for the network itself. It does not require
relative-error approximation of rare probabilities, nor convergence
in total variation to a Gaussian posterior.

\begin{proposition}[The nonpositive-exponent posterior limits]
\label{prop:gp-limits}
For $\gamma=0$, the Boolean posterior converges to that of
$G\sim\mathcal N(0,K)$ conditioned on $|G(x)|>a$ everywhere and on
the observed signs. For $\gamma<0$, it converges to the Gaussian
posterior conditioned only on the observed signs.
\end{proposition}
\begin{proof}
The hidden-kernel argument establishing~\eqref{eq:freeze} holds
for every $\gamma<1$. By~\eqref{eq:readout}, the normalized vector
$N^{\gamma/2}Z_N$ therefore converges in law to
$\mathcal N(0,K)$. Its all-input margin is $aN^{\gamma/2}$, which
equals $a$ when $\gamma=0$ and tends to zero when $\gamma<0$.
The limiting Gaussian is nonsingular. Thus all boundaries involved
in any fixed full-sign event or margin event have probability zero,
and the conditioning events have positive probability. Taking
ratios of their probabilities gives convergence of each Boolean
posterior mass. There are finitely many masses for fixed $d$.
\end{proof}


\section{Fourier level masses of the kernel}\label{app:spectrum}
Put $D=d+1$, $\beta=2^{-(L-1)}$, and $F=\Psi^{\circ(L-1)}$.
The characters $\chi_U(x)=\prod_{j\in U}x_j$ diagonalize the
normalized covariance operator $M^{-1}K$, $M=2^d$. Let $\lambda_j$
be its eigenvalue at degree $j$. We first work on the full $D$-cube
with kernel $F(u\cdot v/D)$, whose degree-$j$ eigenvalue is denoted
by $a_j$. Constants in this and the next section may depend on $L$.
Notation $b\asymp c$ means two inequalities by positive constant
multiples, uniformly in the stated ranges.

\begin{lemma}[Power coefficients and paired eigenvalues]\label{lem:power}
Write $F(z)=\sum_{m\text{ odd}}c_mz^m$. Every odd coefficient is
positive, $\sum_m c_m=1$, and
\[
 c_m\asymp\frac{\Gamma(m-\beta)}{m!}\asymp m^{-1-\beta}
 \qquad(m\ge1\text{ odd}).
\]
The full-cube eigenvalues vanish at even degrees and strictly decrease
through odd degrees. They are positive at every odd degree. On the
original $d$-cube,
\[
 \lambda_{2k}=\lambda_{2k+1}=a_{2k+1}>0,
\]
with a final unpaired eigenvalue if $d$ is even. In particular $K$
is positive definite.
\end{lemma}
\begin{proof}
The series of $\Psi$ has positive coefficients at every odd exponent,
with
\[
 [z^{2k+1}]\Psi(z)=\frac2\pi\frac{\binom{2k}{k}}{4^k(2k+1)}.
\]
Composition preserves strict positivity and oddness. Near $z=1$ in
a slit neighborhood,
\[
 \Psi(z)=1-b(1-z)^{1/2}+O((1-z)^{3/2}),
 \qquad b=\frac{2\sqrt2}{\pi}.
\]
Induction gives
$F(z)=1-A(1-z)^\beta+O((1-z)^{3\beta})$, with $A>0$;
at each composition its leading constant changes from $A_t$ to
$b\sqrt{A_t}$, starting from $A_1=b$.
These expansions admit algebraic coefficient transfer
\cite[Chapter VI]{flajolet2009}. Here are the analytic-domain details.
On the unit circle outside $\{-1,1\}$, positivity of every odd
coefficient gives $|\Psi(z)|<1$: equality in the triangle inequality
would require $z^2=1$. Hence the iterated functions extend
analytically through every compact arc away from these endpoints.
At $1$, restrict $\arg(1-z)$ to a closed proper subinterval of
$(-\pi,\pi)$. Each square-root composition reduces this angular
range, and its positive leading coefficient keeps its image in the
slit neighborhood for the next arcsine composition. Oddness gives
the identical argument at $-1$. These local sectors and the compact
arcs provide the required two-singularity indented domain.
The contributions of $1$ and $-1$ add at odd exponents, giving
\[
 c_m\sim\frac{2A}{-\Gamma(-\beta)}m^{-1-\beta}
 \quad(m\to\infty\text{ odd}).
\]
Together with positivity of the finitely many initial coefficients
and the gamma-ratio asymptotic, this proves the coefficient comparisons.
For $L=2$ the explicit formula and Stirling's bounds suffice instead.

Let $p_m(j)$ be the probability that the simple coordinate-flip walk
on the $D$-cube reaches a specified distance-$j$ vertex in $m$ steps.
Expanding $(u\cdot v/D)^m$ shows that its degree-$j$ Fourier coefficient
is $p_m(j)$. Thus $a_j=\sum_{m\text{ odd}}c_mp_m(j)$.
For the walk recursion,
\begin{align*}
 D[p_{m+1}(j)-p_{m+1}(j+2)]
 ={}&j[p_m(j-1)-p_m(j+1)]\\
 &+(D-j-2)[p_m(j+1)-p_m(j+3)].
\end{align*}
Induction from time zero gives nonincreasing probabilities through
each parity class; terms with zero boundary coefficients are omitted.
The term $m=j$ makes the odd-eigenvalue inequalities strict and proves
positivity. Even-degree eigenvalues vanish by parity.
Restricting the distinguished coordinate to $+1$ gives
$\lambda_q=a_q+a_{q+1}$, proving the claims.
\end{proof}

For odd $r$, the exact combined multiplicity of reduced degrees
$j-1,j$ is $\binom d{j-1}+\binom dj=\binom Dj$. Consequently
\begin{equation}\label{eq:paired}
 D_r=\sum_{\substack{j\le r\\j\text{ odd}}}\binom Dj,
 \quad
 \tau_r=\sum_{\substack{j>r\\j\text{ odd}}}\binom Dj a_j,
 \quad \lambda_{r+1}=a_{r+2}.
\end{equation}
The odd cutoff is essential for these exact identities.

\begin{lemma}[Sharp level and tail estimates]\label{lem:sharp}
Uniformly for odd $1\le j\le D/4$ and odd $1\le r\le D/8$,
\begin{align}
 \binom Dj a_j&\asymp j^{-1-\beta},\label{eq:level}\\
 \tau_r&\asymp r^{-\beta},\label{eq:trace}\\
 \frac{\lambda_{r+1}}{\tau_r}
 &\asymp\frac1{r\binom D{r+2}},\label{eq:ratio}\\
 s_{2,r}:=\sum_{q>r}\binom dq\lambda_q^2
 &\le\frac{C r^{-2\beta-2}}{\binom D{r+2}}.
 \label{eq:secondmoment}
\end{align}
The same level-mass estimate holds up to $j\le\rho D$ for any fixed
$\rho<1/2$, with constants depending on $\rho$.
\end{lemma}
\begin{proof}
The gamma integral and Lemma~\ref{lem:power} give, for odd $j$,
\begin{align}
 a_j&\asymp D^{-\beta}\int_0^\infty
 u^{-1-\beta}p(u)^j(1-p(u))^{D-j}\,du,
 \quad p(u)=\frac{1-e^{-2u}}2,\label{eq:poisson}\\
 &\asymp D^{-\beta}\int_0^{1/2}p^j(1-p)^{D-j}w_\beta(p)\,dp,
 \label{eq:betaweight}
\end{align}
where $w_\beta(p)=[-\frac12\log(1-2p)]^{-1-\beta}/(1-2p)$.
Indeed the gamma integral Poissonizes the walk time, and independent
Poisson coordinate counts give independent coordinate parities with
success probability $p(u)$. Since an odd-distance vertex has zero
probability at even walk times, no extra parity factor occurs.
The weight $w_\beta$ is integrable near $1/2$ and comparable to
$p^{-1-\beta}$ on $[0,b]$ for every fixed $b<1/2$.

Fix $\rho<b<1/2$ and take $j\le\rho D$. The integral on $[0,b]$
is comparable to the complete beta integral
$B(j-\beta,D-j+1)$. For the lower comparison, a beta random variable
with these parameters has mean at most $\rho$, so Markov's inequality
puts at least a fraction $1-\rho/b$ of its mass in $[0,b]$.
The identity
\[
 \binom Dj B(j-\beta,D-j+1)
 =\frac{\Gamma(D+1)\Gamma(j-\beta)}
        {\Gamma(j+1)\Gamma(D+1-\beta)}
 \asymp D^\beta j^{-1-\beta}
\]
gives the desired mass. On $[b,1/2]$, the binomial factor decreases
with $p$ and, after multiplication by $\binom Dj$, is at most
$\exp[-D\operatorname{KL}(\rho\Vert b)]$. The remaining integrable
weight makes this contribution exponentially small. This proves
\eqref{eq:level} and its stated extended range.

Sum \eqref{eq:poisson} over the odd distances above $r$. Dropping the
parity restriction and using the binomial Markov bound gives
\[
 \tau_r\le C D^{-\beta}\int_0^\infty
 u^{-1-\beta}\min\{1,Du/(r+1)\}\,du
 \le C(r+1)^{-\beta}.
\]
For the lower bound, sum \eqref{eq:level} over the odd $j$ in
$(r,2r+2]$, using a fixed cutoff strictly below $D/2$ as needed at
the endpoint. This yields $c r^{-\beta}$; bounded $r$ are absorbed
by a change of constants. Equation~\eqref{eq:ratio} follows from
\eqref{eq:level}, \eqref{eq:trace}, and \eqref{eq:paired}.

For the second moment, use paired multiplicities to write
$s_{2,r}=\sum_{j>r,\ j\text{ odd}}\binom Dj a_j^2$.
Up to an odd cutoff $b_D$ near $D/4$, its terms are at most
$Cj^{-2\beta-2}/\binom Dj$. Consecutive odd inverse binomial
coefficients have ratio
\[
 \frac{(j+1)(j+2)}{(D-j)(D-j-1)},
\]
uniformly less than a fixed constant below one in this range.
The sum is therefore bounded by a constant times its first term,
at $j=r+2$. Beyond $b_D$, monotonicity gives the bound
$a_{b_D}\tau_{b_D}\le C D^{-2\beta-1}/\binom D{b_D}$,
which is exponentially smaller than the claimed bound uniformly
for $r\le D/8$. This proves \eqref{eq:secondmoment}.
\end{proof}

Finally, two-step binomial ratios give, uniformly for odd $r\le D/8$,
\begin{equation}\label{eq:dimensionratios}
 \binom Dr\le D_r\le C\binom Dr,
 \qquad \frac{D_r}{\binom D{r+2}}\le C(r/D)^2.
\end{equation}
All constants used here can be fixed independently of a subsequently
chosen small positive limit $r/d\to\alpha$.


\section{Dictators separate the subcritical regimes}
\label{sec:dictators}
The smallest nonconstant parity already shows that a lower bound
$d^{c k}$ cannot hold uniformly over all $k\ge1$ throughout
$\gamma<1$. Here and below the width tends to infinity at fixed
$d$, training sample, and exponent. We use the normalized Walsh
eigenvalues $\lambda_j$ of $K/2^d$.

\begin{lemma}[Exact ground states]\label{lem:dictator-ground}
The largest normalized eigenvalue is
$\lambda_*:=\lambda_0=\lambda_1$, and its eigenspace consists
exactly of affine functions on the cube. Moreover,
$\lambda_*\le1/(d+1)$. Among all vectors with $|z(x)|\ge a$
at every input, the minimum of $z^{\mathsf T}K^{-1}z/2$ is
$a^2/(2\lambda_*)$. Its minimizers are precisely
\[
 z\in\{\pm a,\ \pm a x_1,\ldots,\pm a x_d\}.
\]
\end{lemma}
\begin{proof}
The pairing and strict spectral ordering follow from
Lemma~\ref{lem:power}. For an elementary verification of strictness,
put $D=d+1$ and write $F(t)=\sum_{m\text{ odd}}c_mt^m$, where
every odd coefficient is positive. The degree-$j$ eigenvalue on the
full $D$-cube is $a_j=\sum_m c_mp_m(j)$, with $p_m(j)$ the
coordinate-flip walk probability of a specified distance-$j$ vertex.
Its exponential generating formula is
\[
 p_m(j)=\frac{m!}{D^m}[t^m](\sinh t)^j(\cosh t)^{D-j}.
\]
Subtracting the formula at $j+2$ leaves
$(\sinh t)^j(\cosh t)^{D-j-2}$, whose coefficients are
nonnegative. The coefficient at $m=j$ is positive. Thus
$a_j>a_{j+2}$ at every available odd pair, and restricting the
distinguished coordinate gives $\lambda_{2j}=\lambda_{2j+1}$.
Consequently the top eigenspace on the original cube has dimension
$d+1$ and consists of degrees zero and one. The normalized covariance
has trace $K(x,x)=1$, so $(d+1)\lambda_*\le1$.

For the uniform population norm,
\[
 z^{\mathsf T}K^{-1}z
 =\sum_U\frac{\widehat z(U)^2}{\lambda_{|U|}}
 \ge\frac{\|z\|_2^2}{\lambda_*}
 \ge\frac{a^2}{\lambda_*}.
\]
Equality requires both $|z(x)|=a$ everywhere and membership in the
top eigenspace. If $P=z/a=c_0+\sum_jc_jx_j$ takes values in
$\{-1,1\}$, then $P^2=1$. Its degree-two coefficients give
$c_ic_j=0$ for $i\ne j$, and its degree-one coefficients give
$c_0c_i=0$. Hence at most one of $c_0,c_1,\ldots,c_d$ is nonzero;
the constant coefficient of $P^2$ makes that coefficient $\pm1$.
The listed functions attain equality.
\end{proof}

\begin{theorem}[Logarithmic samples for a dictator]
\label{thm:dictator-fast}
Fix $L\ge2$, $a>0$, $0<\gamma<1$, and $0<\delta<1$.
For target $f(x)=x_1$, let $S$ consist of independent uniform
training inputs. If
\[
 n\ge\left\lceil\log_2\frac{2d}{\delta}\right\rceil,
\]
then, with probability at least $1-\delta$,
\[
 \lim_{N\to\infty}q_{N,S}(x)=1\qquad\text{for every }x\in\mathcal X.
\]
In particular the strict-$2/3$ confidence-failure fraction is zero
for all sufficiently large widths on this event.
\end{theorem}
\begin{proof}
The target is always a compatible ground-state sign. Its opposite
$-x_1$ is excluded by any training observation. Each of the other
$2d$ constants or signed coordinate functions agrees with the target
at a uniform input with probability $1/2$, and hence on all $n$
independent inputs with probability $2^{-n}$. A union bound shows
that the target is the unique compatible ground-state sign with
probability at least $1-2d\,2^{-n}\ge1-\delta$.

On this event, Lemma~\ref{lem:dictator-ground} makes the target the
unique compatible sign function attaining the global minimum
energy. Proposition~\ref{prop:rate} therefore puts posterior mass
tending to one on this sign function. This proves the assertion
simultaneously at all inputs, since the cube is finite at each width
limit. No formula for weights among tied minimizers is used.
\end{proof}

The contrast with the negative-exponent limit concerns actual
posterior probabilities as well. It does not identify confidence with
the sign of a Gaussian posterior mean.

\begin{theorem}[A universal linear-dimensional obstruction for
negative exponents]\label{thm:negative-linear}
Let $f:\mathcal X\to\{-1,1\}$ be any target and $S$ any training
sample of size $n$, allowing repetitions. Under the Gaussian prior
conditioned on the observed signs, set
\[
 F_\theta(S)=\mathbb P_x\{
       \mathbb P(f(x)G(x)>0\mid f(X_i)G(X_i)>0\ \forall i)
                  \le\theta\}.
\]
For every $1/2<\theta<1$,
\begin{equation}\label{eq:negative-linear}
 F_\theta(S)\ge
 1-\left(2+\frac1{\pi(\theta-1/2)^2}\right)\frac{n}{d+1}.
\end{equation}
Consequently, for every fixed $\gamma<0$, the actual neural-network
posterior with the stipulated global-margin convention satisfies
\begin{equation}\label{eq:negative-linear-width}
 \liminf_{N\to\infty}F_N(S)
 \ge1-\left(2+\frac{400}{9\pi}\right)\frac{n}{d+1}.
\end{equation}
In particular $n=o(d)$ leaves confidence failure $1-o(1)$, even
for a dictator or a constant target.
\end{theorem}
\begin{proof}
Remove repeated training inputs and let $s\le n$ be their number.
Write Gaussian regression on these values as
\[
 G(x)=R(x)+\varepsilon(x),\qquad
 R=K_{\mathcal X,S}K_{SS}^{-1}G_S,
 \qquad\sigma_x^2=\operatorname{Var}\varepsilon(x).
\]
The residual process is independent of $G_S$. Whiten the training
vector by $G_S=K_{SS}^{1/2}\xi$. The observed signs form a
homogeneous cone in the standard Gaussian vector $\xi$. Its radius
is independent of its direction, so conditioning on this cone
preserves $\mathbb E\|\xi\|^2=s$. With $\mathbb E_S^*$ denoting
expectation under that sign conditioning, the regression identity gives
\[
 \|R\|_K^2=\|\xi\|^2,\qquad
 \mathbb E_S^*\|R\|_2^2\le\lambda_*s\le\lambda_*n.
\]
These statements concern the second moment of the random regression
field, not just its posterior mean.

The prior covariance of $R$ is positive semidefinite, bounded by
$K$, and has rank at most $s$. Its normalized trace is at most
$s\lambda_*$. Since $\sigma_x^2=1-\operatorname{Var}R(x)$,
\[
 \mathbb P_x\{\sigma_x^2<1/2\}\le2n\lambda_*.
\]
For every remaining input, conditional Gaussian regression yields
\[
 q_S(x)=\mathbb E_S^*\Phi\!\left(\frac{f(x)R(x)}{\sigma_x}\right),
\]
where $\Phi$ is the standard normal distribution function. The
inequality $|\Phi(t)-1/2|\le |t|/\sqrt{2\pi}$ and
$\sigma_x\ge1/\sqrt2$ imply
\[
 |q_S(x)-1/2|
 \le\frac{\sqrt{\mathbb E_S^*R(x)^2}}{\sqrt\pi}.
\]
Thus $q_S(x)>\theta$ at such an input requires
$\mathbb E_S^*R(x)^2>\pi(\theta-1/2)^2$. Markov's inequality,
integrated over inputs, bounds the fraction of these inputs by
$n\lambda_*/[\pi(\theta-1/2)^2]$. Adding the variance exception
and using $\lambda_*\le1/(d+1)$ proves
\eqref{eq:negative-linear}.

Take $\theta=13/20<2/3$. For $\gamma<0$,
Proposition~\ref{prop:gp-limits} gives convergence of each of the
finitely many correct-label posterior probabilities to the
sign-conditioned Gaussian value. Every input with limiting value
at most $13/20$ therefore has finite-width value below $2/3$ for all
sufficiently large widths. This proves
\eqref{eq:negative-linear-width}. The case $n=0$ follows directly
from sign symmetry and also satisfies the displayed bounds.
\end{proof}

\section{An exact pairing of parity sizes}
\label{sec:projective-symmetry}
The following symmetry holds before taking any width limit.
\begin{proposition}[Projective-cube symmetry]
\label{prop:projective-symmetry}
At every width and exponent for which the conditioning is defined,
the iid-sample learning laws of coordinate parities of sizes
$0$ and $1$, of sizes $2$ and $3$, and so on, coincide whenever
both sizes are available in ambient dimension $d$. The equality
includes the distribution of the actual confidence-failure fraction.
\end{proposition}
\begin{proof}
Fix a coordinate $i$ and define an involution of the cube by
\[
 (T_ix)_i=x_i,\qquad (T_ix)_j=x_ix_j\quad(j\ne i).
\]
If $P_i$ swaps augmented coordinates $0$ and $i$, then
$(1,T_ix)=x_iP_i(1,x)$. For fixed weights $W$, replace only the first
matrix by $W'^1=W^1P_i$. Oddness of every hidden activation and of
the readout gives, inductively,
\[
 z_W(T_ix)=x_i z_{W'}(x).
\]
The laws of $W$ and $W'$ agree because the first-layer columns are
exchangeable. For a target $g$, define $g'(T_ix)=x_ig(x)$.
Absolute margins are invariant under this transformation, and the
training-label event for $(g',T_iS)$ becomes exactly the event for
$(g,S)$. Thus the correct-label posterior probabilities and their
failure fractions agree after the change of variables. The involution
preserves the uniform input distribution.

For $g=\chi_U$, $|U|=k$, writing $u=T_ix$ gives
\[
 g'(u)=u_i^{k+1}\prod_{j\in U\setminus\{i\}}u_j.
\]
If $k$ is odd and $i\in U$, this is a $(k-1)$-coordinate parity.
If $k$ is even and $i\notin U$, it is a $(k+1)$-coordinate parity.
Ordinary coordinate permutations identify all supports of a fixed
size, proving the assertion.
\end{proof}

These results refute a $d^{c k}$ lower bound asserted uniformly
including $k=1$ across the entire subcritical regime. They do not
refute a lower bound restricted to $k\to\infty$, and the symmetry
reduces parity size only within one adjacent even--odd pair.

% Gaussian-posterior part of the exponential sample lower-bound theorem.
% Requires theorem, lemma, proposition, corollary environments and amsmath,
% amssymb, amsthm. The spectral/random-Gram lemma is supplied separately.

\section{The actual Gaussian sign posterior}
Let $\mathcal X=\{-1,1\}^d$, $M=2^d$, and
$f(x)=\prod_{j\in U}x_j, \qquad |U|=k$. All function norms and expectations on
$\mathcal X$ use its uniform probability measure. Let $G$ be a
nondegenerate centered Gaussian vector indexed by $\mathcal X$, with
covariance matrix $K$. The training inputs $X_1,\ldots,X_n$ are independent
and uniform. We condition on
\[
 \mathcal A_S=\{f(X_i)G(X_i)>0:1\le i\le n\}.
\]
The prediction at $x$ is its actual conditional correct-label probability
$q_S(x)=\mathbb P(f(x)G(x)>0\mid\mathcal A_S)$, and we write
\[
 F_\theta(S)=\mathbb P_x\{q_S(x)\le\theta\},\qquad
 \theta=\frac{13}{20},\qquad
 c_\theta=\sqrt{2\pi}\left(\frac1e-\frac7{20}\right)>0.
\]
The threshold $\theta<2/3$ supplies a strict buffer for passage from
limiting Gaussian posteriors to finite-width posterior probabilities.

Suppose the Walsh characters are eigenfunctions of the normalized
covariance operator, with degree-$j$ eigenvalue $\lambda_j$. Write
$V_r=\operatorname{span}\{\chi_U:|U|\le r\}$ and $D_r=\dim V_r$.
For $r<d$, let $Q$ be the covariance contributed by the complementary
Fourier modes, and put
\[
 \tau=Q(x,x)=\sum_{j>r}\binom dj\lambda_j,
 \qquad \ell=\max_{j>r}\lambda_j.
\]
The matrix eigenvalues of $Q$ are $M\lambda_j$; this normalization is
used throughout. Define $\|h\|_K^2=h^TK^{-1}h$, and write $u_+=\max(u,0)$ and $u_-=\max(-u,0)$.

\begin{lemma}[Mean norm and a confidence margin]\label{lem:schedule-mean}
Let $m(x)=\mathbb E[G(x)\mid\mathcal A_S]$, and let $s\le n$ be the
number of distinct sampled inputs. Then
\[
 \|m\|_K^2\le s,\qquad f(X_i)m(X_i)>0.
\]
Let $\sigma_x^2=\operatorname{Var}(G(x)\mid G_S)$, where $S$ here lists
the distinct inputs. Whenever $\sigma_x>0$ and $q_S(x)>\theta$,
\[
 f(x)m(x)>c_\theta\sigma_x.
\]
\end{lemma}
\begin{proof}
Whiten the Gaussian training vector as $G_S=K_{SS}^{1/2}\xi$.
The sign constraints form a homogeneous cone in $\xi$. Conditioning
leaves the squared Gaussian radius with mean $s$, so Jensen's inequality
and Gaussian regression give
\[
 m_S^TK_{SS}^{-1}m_S\le s,\qquad
 m=K_{\mathcal X,S}K_{SS}^{-1}m_S.
\]
The RKHS norm of this regression vector is the preceding quadratic
form. Its observed signed coordinates are strictly positive.

For the confidence claim, the scalar conditional variable
$Y=f(x)G(x)$ has a log-concave density: it is a marginal of a Gaussian
restricted to a convex set, by the marginal log-concavity theorem \cite{prekopa1973}. If $\mu=\mathbb E Y$ and $F_Y$ is its
cdf, log-concavity and Jensen's inequality imply
\[
 \log F_Y(\mu)\ge\mathbb E\log F_Y(Y)=-1.
\]
Consequently $\mathbb P(Y>\mu)\le1-e^{-1}<\theta$.
Conditioning first on the full training vector shows that the density
of $Y$ is a mixture of Gaussian densities with common variance
$\sigma_x^2$, and is therefore at most
$(\sqrt{2\pi}\sigma_x)^{-1}$. If $q_S(x)>\theta$, then $\mu>0$ and
\[
 \theta<\mathbb P(Y>0)
 \le 1-e^{-1}+\frac{\mu}{\sqrt{2\pi}\sigma_x}.
\]
This gives the stated margin.
\end{proof}

\begin{lemma}[Most residual variances remain large]\label{lem:schedule-var}
For every training sample,
\[
 \mathbb P_x\{\sigma_x^2<\tau/2\}\le \frac{2n\ell}{\tau}.
\]
Also, if $p$ is the degree-$r$ Fourier truncation of $m$ and $t=m-p$,
then $\|t\|_2^2\le n\ell$.
\end{lemma}
\begin{proof}
Write $G=G_{\le r}+G_{>r}$ as a sum of independent fields. Conditioning
additionally on the full low-degree field can only decrease conditional
Gaussian variance. The remaining covariance is
\[
 Q-Q_{\mathcal X,S}Q_{SS}^{\dagger}Q_{S,\mathcal X}.
\]
Its subtracted term is positive semidefinite, bounded above by $Q$,
and has rank at most $s\le n$. Its normalized trace is at most $n\ell$.
Markov's inequality proves the first assertion; it includes observed
points automatically. Fourier orthogonality and the mean norm bound
prove $\|t\|_2^2\le\ell\|m\|_K^2\le n\ell$.
\end{proof}

\section{Uniform sampling detects the negative part}
\begin{lemma}[Uniform ReLU sampling on the cube]\label{lem:schedule-relu}
For every fixed $\delta,\eta\in(0,1)$, there is a universal constant
$C$ such that, if
\[
 n\ge C\eta^{-2}D_r\log\!\left(\frac{D_r}{\delta\eta}\right),
\]
then with probability at least $1-\delta$, simultaneously for all
$p\in V_r$,
\[
 \left|\frac1n\sum_{i=1}^n(-f(X_i)p(X_i))_+
             -\mathbb E_x(-f(x)p(x))_+\right|
 \le\eta\mathbb E_x|p(x)|.
\]
In particular, when $r<k$ and $\eta=1/4$, the empirical negative
part is at least $\mathbb E|p|/4$.
\end{lemma}
\begin{proof}
Let $A$ have rows $f(x)(\chi_U(x))_{|U|\le r}$, one for each cube
point. Then $A^TA=MI$ and every row has squared norm $D_r$.
The $\ell_1$ Lewis weights of $A$ are exactly $D_r/M$: substitution
in their defining equation
\[
 w_x^2=A_x(A^T\operatorname{diag}(w^{-1})A)^{-1}A_x^T
\]
verifies this constant solution, and Lewis weights are unique.
The additive ReLU coreset theorem of Mai, Musco, and Rao
\cite[Theorem~5]{mai2021coresets} therefore samples precisely the uniform
distribution on the cube and gives the stated estimate. Its sample bound
is $O(\eta^{-2}D_r\log(D_r/(\delta\eta)))$; no additional logarithmic
factor is required. Finally $\mathbb E fp=0$, because $r<k$, so
$\mathbb E(-fp)_+=\mathbb E|p|/2$.
\end{proof}


\section{A uniform bound through the gaps between Fourier levels}
We prove a lower bound for every sample size using a thin middle band
of Fourier modes whose covariance trace is small. Ordinary Lipschitz
contraction controls this band, while Lewis-weight sampling controls
the lower degrees.

\begin{lemma}[A fixed-threshold sampling bound]\label{lem:multiband-process}
Let $r<u<d$, and put $\Delta=\tau_r-\tau_u$. Let
$\mathcal B_{r,u}$ be the span of Walsh characters of degrees $r+1$
through $u$. Suppose
\[
 \eta^2\ge C\frac{D_r\log(2n)}n
\]
for a sufficiently large universal constant $C$. For deterministic
$R,T\ge0$, the class
\[
 \mathcal H=\{-f(p+b):p\in V_r,\ \mathbb E|p|\le R,
                 \ b\in\mathcal B_{r,u},\ \|b\|_K\le\sqrt n\}
\]
satisfies
\begin{equation}\label{eq:multiband-process}
 \mathbb E_S\sup_{h\in\mathcal H}
   |\mathbb P_n(h-T)_+-\mathbb E_x(h-T)_+|
 \le4(R\eta+\sqrt\Delta).
\end{equation}
Here $\mathbb P_n g=n^{-1}\sum_i g(X_i)$.
\end{lemma}
\begin{proof}
For the full-cube matrix with rows $f(x)(\chi_U(x))_{|U|\le r}$,
the Lewis weights are $D_r/M$. The linear-process moment bound of
Mai, Musco, and Rao \cite[Theorem 3]{mai2021coresets}, with its failure
parameter set to $1/2$, therefore implies
\[
 \mathbb E_{S,\epsilon}\sup_{p\in V_r,\,\mathbb E|p|\le R}
 \left|\frac1n\sum_i\epsilon_i f(X_i)p(X_i)\right|\le R\eta.
\]
The normalization follows directly: their sampled matrix rescales each
sampled row by $M/n$, whereas $\|p\|_1=M^{-1}\|A\beta\|_1$.
Their moment bound implies the displayed expectation bound by Jensen.

The covariance-weighted feature vector of the band has squared norm
$\Delta$ at every input. Consequently
\[
 \mathbb E_{S,\epsilon}\sup_{\|b\|_K\le\sqrt n}
 \left|\frac1n\sum_i\epsilon_i f(X_i)b(X_i)\right|
 \le\frac{\sqrt n}{n}\sqrt{n\Delta}=\sqrt\Delta.
\]
The map $z\mapsto(z-T)_+$ is $1$-Lipschitz and vanishes at zero.
Symmetrization, followed by the scalar contraction principle, bounds
the left-hand side of \eqref{eq:multiband-process} by four times the
linear Rademacher width of $\mathcal H$. The triangle inequality and
the preceding two estimates prove the claim. Both parameter balls and
the threshold are fixed before sampling; no positive homogeneity in
$T$ or relative bound on an unbounded parameter class is asserted.
\end{proof}


\section{A polynomial-threshold obstruction}
The following elementary polynomial-threshold obstruction is due to
Aspnes, Beigel, Furst, and Rudich~\cite{ABFR}; we include its proof.
\begin{lemma}[Polynomial thresholds cannot approximate parity]
\label{lem:ptf-parity}
On a cube of dimension $k$, for $0\le s<k$, every degree-$s$ polynomial threshold, with a fixed common convention
for zero ties, has full-parity error at least
\begin{equation}\label{eq:ptf-parity}
 2^{-k}\sum_{j=0}^{\lfloor(k-s-1)/2\rfloor}\binom kj.
\end{equation}
This lower bound tends to $1/2$ whenever $s=o(\sqrt k)$.
\end{lemma}
\begin{proof}
A sufficiently small constant perturbation resolves any zero ties
without changing the specified threshold classifier. Thus suppose
$P$ has no zeros. Put $t=\lfloor(k-s-1)/2\rfloor$. If its error set
had fewer than $\sum_{j\le t}\binom kj$ points, a nonzero polynomial
$Q$ of degree at most $t$ could be chosen to vanish on every error.
The function $fPQ^2$ would then be nonnegative everywhere and positive
somewhere. But $PQ^2$, after multilinear reduction, has degree at most
$s+2t<k$, so its parity Fourier coefficient
$\mathbb E fPQ^2$ is zero, a contradiction. This proves
\eqref{eq:ptf-parity}. The last assertion follows from the central
binomial bound $\max_j2^{-k}\binom kj=O(k^{-1/2})$.
\end{proof}


\section{Random tail covariance matrices}\label{app:gram}
Let $Q$ be the covariance kernel retaining only degrees greater than an
odd cutoff $r$, and write $\tau=\tau_r$, $\ell=\lambda_{r+1}$, and
$s_2=s_{2,r}$. This section proves the sample event needed later, without
assuming that a Gaussian approximation remains valid after conditioning.
The entire argument concerns the exact limiting Gaussian covariance.

\begin{lemma}[Far-pair decay]\label{lem:far}
There are constants $c_0,C>0$ such that, when $r\le D/8$ is odd,
\[
 \max_{d/4\le\dist(x,y)\le3d/4}|Q(x,y)|
 \le\rho_r:=C\sqrt D\,e^{-r/8}+Ce^{-c_0D}.
\]
\end{lemma}
\begin{proof}
The full-cube Krawtchouk polynomial satisfies
\[
 \sum_{j=0}^D\mathsf K_j(h;D)z^j=(1+z)^{D-h}(1-z)^h.
\]
If $\delta D\le h\le(1-\delta)D$, put $b=\min(h,D-h)$.
On $|z|=t$ the generating function is bounded in absolute value by
\[
 (1+t^2)^b(1+t)^{D-2b}
 \le(1+t)^D\exp\!\left[-\frac{2\delta Dt}{(1+t)^2}\right].
\]
Cauchy's coefficient bound at $t=j/(D-j)$ and the uniform Stirling
lower bound for $\binom Dj$ therefore give, for $1\le j<D$,
\begin{equation}\label{eq:kraw}
 \frac{|\mathsf K_j(h;D)|}{\binom Dj}
 \le C\sqrt D\exp[-2\delta j(1-j/D)].
\end{equation}
For reduced-cube points at distance in $[d/4,3d/4]$, their augmented
full-cube distance lies in $[D/8,7D/8]$ for large $d$.
Because $r$ is odd, the exact tail is
$Q(x,y)=\sum_{j>r,\ j\text{ odd}}a_j\mathsf K_j(h;D)$.
For $r<j\le D-r$, \eqref{eq:kraw} gives a bound
$C\sqrt D e^{-r/8}$ after summation, since the total spectral mass
is one. For $j>D-r$, monotonicity and an odd cutoff $b_D$ near $D/4$
give
\[
 \sum_{j>D-r}\binom Dj a_j
 \le a_{b_D}\sum_{k<r}\binom Dk
 \le C D^{-\beta-1/2}
       \exp[-D(\mathsf h(1/4)-\mathsf h(1/8))+O(1)].
\]
Here $\mathsf h(t)=-t\log t-(1-t)\log(1-t)$ denotes the binary entropy in nats.
The positive entropy gap absorbs the polynomial factor and proves the
claim. Adjust constants for bounded $d$.
\end{proof}

\begin{lemma}[A tail Gram deviation bound]\label{lem:gram-general}
Let $X_1,\ldots,X_n$ be independent uniform cube points, with
$1\ll n\le M/4$. Let
\[
 B_{\mathrm{near}}=\#\{y:\dist(x,y)\notin[d/4,3d/4]\}
 \le 2e^{d\mathsf h(1/4)},\qquad q=2\log M.
\]
The count is independent of $x$. For every $z>0$,
\begin{align}
 &\mathbb P\{\|Q_{SS}-\tau I_n\|_{\mathrm{op}}>z\tau\}
 \notag\\
 &\quad\le\frac C{z^2}
 \left[\frac{q\rho_r}{\tau}
       +\frac{\sqrt{nq s_2}}{\tau}
       +\frac{n\ell}{\tau}
       +\frac{n(1+B_{\mathrm{near}})}M\right]^2
       +\frac{n^2(1+B_{\mathrm{near}})}M.
 \label{eq:gram-bound}
\end{align}
Here the sampled matrix includes repeated training locations.
\end{lemma}
\begin{proof}
Let $J=Q-\tau I_M$, a zero-diagonal Hermitian matrix. Form a deterministic
mask $J_{\mathrm{far}}$ retaining only pairs at distances in
$[d/4,3d/4]$. Since $Q$ is positive semidefinite with constant diagonal
$\tau$, every entry has absolute value at most $\tau$. Thus
\begin{align*}
 \|J_{\mathrm{far}}\|_{\max}&\le\rho_r,\\
 \|J_{\mathrm{far}}\|_{1\to2}^2&\le M s_2,\\
 \|J_{\mathrm{far}}\|_{\mathrm{op}}
 &\le M\ell+\tau+\tau B_{\mathrm{near}}.
\end{align*}
For the second inequality, Parseval gives squared column norm
$M s_2-\tau^2$ for $J$, and entrywise masking only decreases it.
For the third, $\|J\|\le M\ell+\tau$, while the removed matrix has
maximum absolute row sum at most $\tau B_{\mathrm{near}}$.
The entropy bound on $B_{\mathrm{near}}$ is the elementary binomial
tail bound on distances to a fixed cube point.

Let $R$ retain each of the $M$ coordinates independently with
probability $p=2n/M$. The zero-diagonal case of Tropp's random
principal-submatrix theorem \cite[Theorem 1.1]{tropp2008}, at moment
order two and $q=2\log M$, states
\[
 (\mathbb E\|RJR\|^2)^{1/2}
 \le C\left[q(\mathbb E\|RJR\|_{\max}^2)^{1/2}
       +\sqrt{pq}(\mathbb E\|JR\|_{1\to2}^2)^{1/2}
       +p\|J\|\right].
\]
Apply it to $J_{\mathrm{far}}$ and substitute the preceding bounds:
\begin{equation}\label{eq:bernoulli-gram}
 (\mathbb E\|R J_{\mathrm{far}}R\|^2)^{1/2}
 \le C[q\rho_r+\sqrt{nq s_2}+n\ell
               +(n/M)\tau(1+B_{\mathrm{near}})].
\end{equation}
The logarithm here is that of the original dimension $M$, not the
sample size.

To pass to a uniform subset of size $n$, use a uniform permutation of
all cube points, independently of $B\sim\operatorname{Bin}(M,2n/M)$.
Its first $B$ elements form the Bernoulli subset. On $B\ge n$, the
first $n$ principal submatrix is a compression of the larger one, so
its operator norm cannot be larger. The size-$n$ event is independent
of $B$, and $\mathbb P(B\ge n)\ge1/2$ for large $n$. Thus Markov's
inequality applied to \eqref{eq:bernoulli-gram} gives the same fixed-size
probability bound with a constant factor loss.

Conditional on no collisions, an independent size-$n$ sample is a
uniform size-$n$ subset, up to an irrelevant ordering. Collisions have
probability at most $n(n-1)/(2M)$. The probability of any sampled pair
violating the distance condition is at most $n^2B_{\mathrm{near}}/M$.
On their complement, the sampled matrices $J$ and $J_{\mathrm{far}}$
agree. Adding these exceptional probabilities proves
\eqref{eq:gram-bound}.
\end{proof}


\section{Tail covariance estimates at sparse cutoffs}\label{sec:sparse-gram}
The original principal-submatrix estimate uses the logarithm of the full
cube cardinality. We retain that estimate here. A more precise bound on
entries near orthogonal input pairs makes it applicable even when the
spectral cutoff is a sufficiently large fixed integer.

\begin{lemma}[Near-orthogonal entries]\label{lem:near-orthogonal}
There are constants $a_0,\delta_0,c,C>0$ such that the following holds.
Put $D=d+1$, let $r$ be odd with $r+2\le\delta_0D$, and let $h$ be the
Hamming distance between the augmented inputs $(1,x),(1,y)$.
If $a=|D-2h|/D\le a_0$, then
\begin{equation}\label{eq:near-orthogonal}
 \frac{|Q_r(x,y)|}{\tau_r}
 \le [C(a+\sqrt{(r+2)/D})]^{r+2}+C d^\beta e^{-cd},
 \qquad Q_r=K_{>r}.
\end{equation}
\end{lemma}
\begin{proof}
For $1\le j<D$, the Krawtchouk generating function gives
\begin{equation}\label{eq:sparse-kraw}
 \frac{|\mathsf K_j(h;D)|}{\binom Dj}
 \le [e^{3/2}(a+\sqrt{j/D})]^j.
\end{equation}
Indeed, set $b=\min(h,D-h)$. On the circle $|z|=t$, the generating
function has modulus at most
\[
 (1+t^2)^b(1+t)^{D-2b}\le \exp(Dt^2/2+Dat).
\]
Cauchy's coefficient bound and $\binom Dj\ge(D/j)^j$, with
$t=j/[D(a+\sqrt{j/D})]$, prove \eqref{eq:sparse-kraw}: both $Dt^2$
and $Dat$ are at most $j$.

Choose $a_0,\delta_0$ so small that
$e^{3/2}(a_0+\sqrt{\delta_0})\le e^{-1}$ and $\delta_0<1/8$.
On $j\le\delta_0D$, the derivative of
$j\log[e^{3/2}(a+\sqrt{j/D})]$ is at most $-1+1/2<0$.
Consequently \eqref{eq:sparse-kraw} is bounded by its value at $r+2$
throughout $r<j\le\delta_0D$. The sum of the corresponding positive
spectral masses $\binom Dj a_j$ is at most $\tau_r$.
For $\delta_0D<j\le(1-\delta_0)D$, the central-distance estimate
\eqref{eq:kraw} bounds the normalized Krawtchouk polynomial by
$C\sqrt D e^{-cD}$; its polynomial factor can be absorbed by decreasing
$c$. For $j>(1-\delta_0)D$, choose an odd $b_D$ near $D/4$.
Monotonicity of $a_j$ and Lemma~\ref{lem:sharp} bound the remaining
mass by
\[
 a_{b_D}\sum_{i<\delta_0D}\binom Di
 \le \operatorname{poly}(D)
       e^{-D[\mathsf h(1/4)-\mathsf h(\delta_0)]},
\]
where $\mathsf h(t)=-t\log t-(1-t)\log(1-t)$.
The entropy gap is positive. Finally $\tau_r\ge cD^{-\beta}$ by
Lemma~\ref{lem:sharp}, proving \eqref{eq:near-orthogonal}.
\end{proof}

\begin{lemma}[Sparse adaptive tail Gram matrices]\label{lem:sparse-gram}
Fix $A\ge1$, a failure probability $\delta>0$, and an odd integer
$r_0\ge11$. There is $\alpha_0>0$ such that, for any fixed
$0<\alpha\le\alpha_0$, the following assertions hold uniformly as
$d\to\infty$. Suppose $r$ is odd,
\[
 r_0\le r\le\alpha d,\qquad
 A D_r\log(D_r+1)\le n<A D_{r+2}\log(D_{r+2}+1).
\]
Put $u=r+2$, $\rho=r/(d+1)$, and $H=\log(e/\rho)$.
For each such specified $n$, with probability at least $1-\delta$,
\begin{align}
 \|(Q_r)_{SS}\|_{\mathrm{op}}
 &\le C_{L,\delta}\left[1+\frac{\sqrt{nd\,s_{2,r}}}{\tau_r}
                   +\frac{n\ell_r}{\tau_r}\right]\tau_r
 \notag\\
 &\le C_{L,\delta}
       [1+\sqrt{A\rho^{-1}H}+AH]\tau_r,
       \label{eq:sparse-coarse-gram}\\
 \|(Q_u)_{SS}-\tau_uI_n\|_{\mathrm{op}}
 &\le\bigl[C_{L,\delta}
       (\sqrt{A\rho H}+A\rho^2H)+o(1)\bigr]\tau_r.
       \label{eq:sparse-fine-gram}
\end{align}
The $o(1)$ is uniform over the stated range. In particular, after
reducing $\alpha_0$ if necessary,
$\|(Q_u)_{SS}\|_{\mathrm{op}}\le2\tau_r$ on the same event for all
sufficiently large $d$. Deterministically,
\begin{align}
 0\le\tau_r-\tau_u&\le C_L\tau_r/r,
       \label{eq:gram-sparse-band-trace}\\
 n\ell_u/\tau_r&\le C_LA\rho^2H,
       \label{eq:sparse-fine-eigenvalue}
\end{align}
where $\ell_v=\lambda_{v+1}$. No simultaneous event over different
sample sizes is asserted.
\end{lemma}
\begin{proof}
The dimension and spectral estimates already proved give
\begin{align*}
 n&\le C A\binom D{r+2}rH,&
 \log(2n)&\le C_A rH,\\
 n\ell_r/\tau_r&\le C_LAH,&
 nd\,s_{2,r}/\tau_r^2&\le C_LA\rho^{-1}H.
\end{align*}
For $u=r+2$, the exact two-step binomial ratio gives
\begin{align*}
 n\ell_u/\tau_r&\le C_LA\rho^2H,&
 nd\,s_{2,u}/\tau_r^2&\le C_LA\rho H.
\end{align*}
Moreover $\tau_r-\tau_u=\binom D{r+2}a_{r+2}$, so the level and
trace estimates prove \eqref{eq:gram-sparse-band-trace}.

We apply the proof of Lemma~\ref{lem:gram-general} with a different
deterministic mask. Retain pairs satisfying
\[
 \left|\frac{1+x\cdot y}{d+1}\right|\le a_n,
 \qquad a_n=\sqrt{128\log(2n)/d}.
\]
By choosing $\alpha_0$ sufficiently small, uniformly
$\log(2n)/d\le C_A\alpha\log(e/\alpha)+o(1)$ is as small as needed,
so $a_n\le a_0$. For independent uniform inputs, Hoeffding's bound
(and $a_n\ge2/(d+1)$) gives exceptional-pair probability
\[
 b_n\le2\exp(-a_n^2d/8)=2(2n)^{-16}.
\]
The exceptional set is translation invariant on the cube. Its size
is therefore the same in every row, so the masked-matrix proof applies
with $B_{\mathrm{near}}/M=b_n$. This mask is deterministic for each
fixed pair $(d,n)$; it does not depend on the observed sample.
Lemma~\ref{lem:near-orthogonal} supplies the retained-entry bound
\begin{equation}\label{eq:sparse-coherence}
 \frac{\eta_v}{\tau_v}
 \le\left[C\sqrt{\frac{\log(2n)+v+2}{d}}\right]^{v+2}
       +C d^\beta e^{-cd},\qquad v\in\{r,u\}.
\end{equation}
Thus exactly the argument of Lemma~\ref{lem:gram-general}, still with
$q=2\log M=2d\log2$, gives for every $z>0$
\begin{align}
 &\mathbb P\{\|(Q_v)_{SS}-\tau_vI_n\|>z\tau_v\}\notag\\
 &\quad\le\frac C{z^2}
 \left[\frac{q\eta_v}{\tau_v}
       +\frac{\sqrt{nq s_{2,v}}}{\tau_v}
       +\frac{n\ell_v}{\tau_v}+n(M^{-1}+b_n)\right]^2
       +n^2(M^{-1}+b_n).
 \label{eq:sparse-masked-tropp}
\end{align}
The collision and exceptional-pair terms tend to zero uniformly:
$n\ge A D_{r_0}\log(D_{r_0}+1)\to\infty$, and $\alpha_0$ can be chosen
so that $n\le e^{c_1d}$ with $2c_1<\log2$.

The coherence term $q\eta_v/\tau_v$ also tends to zero uniformly.
For $r\le\sqrt d$, its bracket in \eqref{eq:sparse-coherence} is at
most $C_A d^{-1/4}\sqrt{\log d}$, and its exponent is at least
$r_0+2\ge13$. Hence it is bounded, after multiplication by $q$, by
\[
 C d[C_A d^{-1/4}\sqrt{\log d}]^{13}
       +C d^{1+\beta}e^{-cd}=o(1).
\]
For $r>\sqrt d$, reducing $\alpha_0$ makes the bracket at most $1/2$,
so the bound is $Cd\,2^{-\sqrt d}+Cd^{1+\beta}e^{-cd}=o(1)$.
Using \eqref{eq:sparse-masked-tropp} for $v=r,u$, with failure
probability at most $\delta/2$ for each, now proves
\eqref{eq:sparse-coarse-gram} and \eqref{eq:sparse-fine-gram}.
For the latter, multiply by $\tau_u/\tau_r\le1$ before using the
spectral bounds, so no lower comparison of the two traces is needed.
Finally $\rho H\to0$ as $\rho\to0$, uniformly for $0<\rho\le\alpha$;
choose $\alpha_0$ to make the fixed part of the fine deviation below
$1/2$, and absorb the uniform $o(1)$. Since $\tau_u\le\tau_r$, this
gives the asserted upper bound $2\tau_r$.
\end{proof}

\section{A lower bound for sparse parity at every sample size}
\label{sec:sparse-gaussian}
Let $U\subseteq[d]$, $|U|=k$, and $f=\chi_U$. Write
$F_\theta(S)=\mathbb P_x\{q_S(x)\le\theta\}$ for the actual
Gaussian sign-posterior confidence failure, with
\[
 \theta=\frac{13}{20},\qquad
 c_\theta=\sqrt{2\pi}\left(\frac1e-\frac7{20}\right),\qquad
 a_\theta=\frac{c_\theta}{2\sqrt2}.
\]
The next theorem has a failure fraction independent of both $k$ and $d$.

\begin{theorem}[Ambient-dimension lower bound]
\label{thm:sparse-gaussian}
For each fixed depth $L\ge2$, there are constants $c>0$ and $d_0$
such that, whenever $d\ge d_0$, $1\le k\le d$, and
\[
 0\le n\le\exp\!\left(c k\log\frac{ed}{k}\right),
\]
independent uniform training inputs satisfy
\[
 \mathbb P_S\{F_{13/20}(S)\ge\varepsilon\}\ge\frac34,
 \qquad \varepsilon=\frac{c_\theta^2}{256}>0.
\]
The assertion holds for each specified sample size. In the network model
of the width-limit proposition, every fixed $\gamma<0$ consequently
satisfies the same lower bound for
$\liminf_{N\to\infty}\mathbb P_x\{q_{N,S}(x)\le2/3\}$.
\end{theorem}

\begin{proof}
Fix a failure probability $\delta=1/40$. We choose a sampling constant
$A$ sufficiently large, a degree fraction $\alpha>0$ sufficiently small,
and an odd integer $r_0\ge11$ sufficiently large, in that order.
All these choices depend only on the fixed depth and universal constants.
Further decreases in $\alpha$ are harmless.

We first treat $k\ge K_0$, where the fixed constant $K_0$ is chosen
below; bounded $k$ will be covered at the end.
For a specified $n>0$, let $r$ be the largest odd integer satisfying
\begin{equation}\label{eq:sparse-coarse-cutoff}
 A D_r\log(D_r+1)\le n.
\end{equation}
First suppose this cutoff exists and $r\ge r_0$. By choosing $c$ in the
theorem sufficiently small and then $K_0$ sufficiently large, we ensure
\begin{equation}\label{eq:sparse-cutoff-range}
 r\le\alpha k\le\alpha d,\qquad u:=r+2<k,
 \qquad n<A D_{r+2}\log(D_{r+2}+1).
\end{equation}
Indeed, the binomial lower bound
$\binom dr\ge(d/r)^r$ shows that, for $r\ge\alpha k$, the logarithm
of the sampling threshold is at least a constant multiple of
$\alpha k\log(ed/k)$. Rounding by two can be absorbed by increasing
$K_0$. Also $k\log(ed/k)\le d$, so the same choice of $c$ puts $n$
within the small-exponential range needed for the random Gram bounds.
Uniformly in this shell interval, for large $d$,
\[
 \log(2n)\le C\log(D_r+1)\le C r\log(eD/r),\qquad D=d+1.
\]

Put $\rho=r/D$, $H_\rho=\log(e/\rho)$, $\tau=\tau_r$,
$\ell_r=\lambda_{r+1}$, and
\[
 B=C_{L,\delta}\left(1+
             \frac{\sqrt{nd\,s_{2,r}}}{\tau}
             +\frac{n\ell_r}{\tau}\right),
 \qquad \eta^2=C_0\frac{D_r\log(2n)}n.
\]
Choose $C_0$ as required by Lemma~\ref{lem:multiband-process}.
Lemma~\ref{lem:sparse-gram} and its near-orthogonal coherence estimate
imply, with probability at least $1-\delta$ for all sufficiently large
$d$,
\begin{equation}\label{eq:sparse-low-gram}
 \|Q_{r,SS}\|_{\mathrm{op}}\le B\tau.
\end{equation}
This holds uniformly over $r\ge r_0$ in
\eqref{eq:sparse-cutoff-range}; it does not require $r\gg\log d$.
The sharp spectral ratios give the crucial cancellation
\begin{equation}\label{eq:sparse-variable-accuracy}
 \eta^2B\le C_{L,\delta}\left[
      \frac1A+\sqrt{\frac{\rho H_\rho}{A}}+
                     \rho^2H_\rho\right].
\end{equation}
For completeness, the constant term uses $\eta^2\le C/A$.
The eigenvalue term is bounded by
\[
 \eta^2\frac{n\ell_r}{\tau}
 \le\frac{C D_r\log(2n)}{r\binom D{r+2}}
 \le C\rho^2H_\rho.
\]
The square of the variance term is at most
\[
 \left(\eta^2\frac{\sqrt{nd\,s_{2,r}}}{\tau}\right)^2
 \le\frac CA\frac{C D_r\log(2n)d}{r^2\binom D{r+2}}
 \le\frac CA\rho H_\rho.
\]
Thus $\eta\sqrt B$ can be made arbitrarily small uniformly in $n$
by choosing $A$ large and then $\alpha$ small.

Let $m=\mathbb E[G\mid\mathcal A_S]$ and $p=\Pi_{\le r}m$.
By Lemma~\ref{lem:schedule-mean}, $\|m\|_K^2\le n$ and its training
signs are correct. The ordinary first-power ReLU sampling event in
Lemma~\ref{lem:schedule-relu}, at accuracy $1/4$, has probability at
least $1-\delta$ after increasing $A$. Since $r<k$, Walsh orthogonality
gives $\mathbb E fp=0$, so on this event and
\eqref{eq:sparse-low-gram},
\begin{equation}\label{eq:sparse-head-radius}
 \frac14\mathbb E|p|
 \le\mathbb P_n(-fp)_+
 \le\sqrt{\mathbb P_n(m-p)^2}
 \le\sqrt{B\tau}.
\end{equation}
Set the deterministic radius $R=4\sqrt{B\tau}$.

\medskip\noindent
\emph{Only one additional paired level is required.}
Write
\[
 m=p+b+t,\qquad b=\Pi_{(r,u]}m,\quad t=\Pi_{>u}m,\quad u=r+2,
 \qquad \Delta=\tau_r-\tau_u.
\]
The level-mass estimate and the two-step binomial ratio give
\begin{align}
 \Delta/\tau&\le C_L/r,\label{eq:sparse-band-trace}\\
 n\ell_u/\tau&\le C_L A\rho^2H_\rho,
       \label{eq:sparse-high-eigenvalue}\\
 nd\,s_{2,u}/\tau^2&\le C_L A\rho H_\rho.
       \label{eq:sparse-high-variance}
\end{align}
For example, maximality yields
$n\le C A\binom D{r+2}rH_\rho$, whereas the two high-tail bounds
involve $\binom D{r+4}$ in the denominator; its ratio with
$\binom D{r+2}$ supplies a factor $O(\rho^2)$.
By increasing $r_0$ we have $\tau_u\ge3\tau/4$.
By reducing $\alpha$, Lemma~\ref{lem:sparse-gram} applied at $u$
now gives, with probability at least $1-\delta$,
\begin{equation}\label{eq:sparse-high-gram}
 \|Q_{u,SS}\|_{\mathrm{op}}\le2\tau.
\end{equation}
The coherence remainder tends to zero uniformly, including at fixed
$r_0$. The two displayed high-tail terms can be made as small as needed
before taking $d$ large.

Put $\zeta=n\ell_u/\tau$. We choose $\alpha$ small enough that
\begin{equation}\label{eq:sparse-zeta-choice}
 e_*:=(4+a_\theta^{-2})\zeta\le a_\theta^2/32.
\end{equation}
This is possible uniformly by \eqref{eq:sparse-high-eigenvalue}.
The RKHS bound, the high-tail Gram event, and the variance rank bound
imply
\[
 \mathbb P_n t^2\le2\tau,\qquad
 \|t\|_2^2\le\zeta\tau,
 \qquad
 \mathbb P_x\{\sigma_x^2<\tau/2\}\le4\zeta.
\]
For the last assertion, condition additionally on the full degree-$u$
field. Its remaining diagonal before training is
$\tau_u\ge3\tau/4$, and the normalized trace of the variance reduction
is at most $n\ell_u=\zeta\tau$.
Consequently, except on an input set of measure at most $e_*$, both
$\sigma_x^2\ge\tau/2$ and $|t(x)|\le a_\theta\sqrt\tau$ hold.
The confidence lemma then shows that
\begin{equation}\label{eq:sparse-confidence-margin}
 f(x)(p(x)+b(x))>a_\theta\sqrt\tau
\end{equation}
at every remaining point where $q_S(x)>\theta$.

\medskip\noindent
\emph{A fixed stop-loss bound rules out hiding the negative mass.}
Choose the deterministic threshold
$T=4\sqrt\tau/a_\theta$.
Apply Lemma~\ref{lem:multiband-process} to the deterministic radius
$R=4\sqrt{B\tau}$ and this $T$.
By Markov's inequality its uniform empirical error is at most
\[
 4\delta^{-1}(R\eta+\sqrt\Delta)
 \le a_\theta\sqrt\tau/8
\]
with probability at least $1-\delta$, after the stated choices:
\eqref{eq:sparse-variable-accuracy} controls the first term and
\eqref{eq:sparse-band-trace} controls the second by increasing $r_0$.
The actual $p,b$ belong to this class on the earlier events, since
\eqref{eq:sparse-head-radius} bounds $\mathbb E|p|$ and
$\|b\|_K\le\sqrt n$.
Let $h=p+b$. Training consistency gives $(-f_i h_i)_+\le|t_i|$.
Using $(z-T)_+\le z^2/(4T)$ for $z\ge0$ and
\eqref{eq:sparse-high-gram}, we obtain
\[
 \mathbb P_n(-fh-T)_+
 \le\frac{\mathbb P_n t^2}{4T}
 \le a_\theta\sqrt\tau/8,
 \qquad
 \mathbb E_x(-fh-T)_+\le a_\theta\sqrt\tau/4.
\]
By \eqref{eq:sparse-confidence-margin}, the positive-part integral of
$fh$ is at least
$a_\theta\sqrt\tau(1-F_\theta-e_*)$, and its negative set has measure
at most $F_\theta+e_*$. Crucially, $\deg h\le u<k$, so
$\mathbb E fh=0$ exactly. Equality of the positive and negative
integrals therefore gives
\[
 a_\theta(1-F_\theta-e_*)
 \le\frac4{a_\theta}(F_\theta+e_*)+\frac{a_\theta}4.
\]
It follows that
\[
 F_\theta\ge\frac{3a_\theta^2}{4(4+a_\theta^2)}-e_*
             \ge\frac{a_\theta^2}{32}=\varepsilon,
\]
where we used \eqref{eq:sparse-zeta-choice} and $a_\theta^2<1$.
The low Gram, high Gram, ReLU, and fixed-stop-loss events have combined
probability at least $1-4\delta>3/4$ for all sufficiently large $d$.
No independence between these events is assumed.

\medskip\noindent
\emph{The remaining, smaller samples.}
If the maximal cutoff is below $r_0$, or does not exist, then
$n<A D_{r_0}\log(D_{r_0}+1)$. Fix the odd integer $s=r_0+2$.
The sharp spectrum gives, uniformly in this entire range,
\[
 \frac{n\ell_s}{\tau_s}
 \le C_{L,A,r_0}\frac{D_{r_0}\log(D_{r_0}+1)}
                           {s\binom D{s+2}}=o(1)
 \qquad(d\to\infty).
\]
Let $p=\Pi_{\le s}m$. The mean-norm, residual-variance, and confidence
lemmas show that, except on an $o(1)$ fraction of inputs, confidence
above $\theta$ implies $fp>0$. Conditional on any values of the
$d-k$ nuisance coordinates, $p$ is a degree-$s$ polynomial on the
$k$ target coordinates. Lemma~\ref{lem:ptf-parity} thus bounds its
strictly correct fraction by $1-b_k(s)$, where
\[
 b_k(s)=2^{-k}\sum_{j=0}^{\lfloor(k-s-1)/2\rfloor}\binom kj.
\]
Choose $K_0$ sufficiently large relative to the fixed $s$ that
$b_k(s)\ge2/5$ for every $k\ge K_0$.
This gives $F_\theta\ge2/5-o(1)>\varepsilon$ for every such training
sample. At $n=0$, $q_S=1/2$ identically.
All errors above are uniform over $k\ge K_0$ and the specified sample
range after fixing the constants.

Finally shrink $c$, if necessary, so that $c\le1/(2K_0)$.
For $1\le k<K_0$ the sample bound implies $n\le\sqrt{ed}$.
The degree-zero projection $p_0=\Pi_{\le0}m$ is constant.
The first paired eigenvalues satisfy $\lambda_0=\lambda_1$;
since the total trace is one, $(d+1)\lambda_0\le1$.
Consequently its tail satisfies the exact bound
$\ell_0/\tau_0=\lambda_1/(1-\lambda_0)\le1/d$.
The same mean and residual-variance estimates show that, outside a
fraction $O_\theta(n/d)=o(1)$ of inputs, confidence above $\theta$
implies $fp_0>0$. A constant can be strictly correct on at most half
the cube for any nonconstant parity. Therefore
$F_\theta\ge1/2-o(1)>\varepsilon$ for every training sample in this
remaining range. A single sufficiently large $d_0$ covers all cases.

Finally, the Gaussian sign-posterior width limit holds at fixed
$d,k,S$ for every fixed $\gamma<0$. At any point with
$q_S(x)\le13/20$, finite-width correct-label probabilities are eventually
below $2/3$. Since the cube is finite, the asserted $\liminf$ transfer
holds simultaneously at all these points.
\end{proof}


\section{Scope of the conclusions}
Theorem~\ref{thm:sparse-gaussian}, together with
Proposition~\ref{prop:gp-limits}, proves Theorem~\ref{thm:main}.
Theorems~\ref{thm:dictator-fast} and~\ref{thm:negative-linear} prove
Theorem~\ref{thm:separation}.
Parity has a Boolean circuit with $O(k)$ gates, whereas the negative
exponent posterior requires the stated number of examples to make its
confidence failure vanish. This is a property of the specified posterior;
parity can be identified from $O(d)$ samples by linear algebra over
the two-element field. Searching the supports of size $k$ instead
requires only $O(k\log(ed/k))$ samples at fixed success probability,
since each incorrect parity agrees with the target independently with
probability $1/2$ on each example.

Two questions remain outside the result. At $\gamma=0$, conditioning
retains a fixed margin at all $2^d$ inputs, so the homogeneous-cone
argument used here is unavailable. For $0<\gamma<1$, the posterior
selects energy minimizers, and the dictator example shows that its
learning behavior need not match the sign-conditioned Gaussian regime.
Neither the proof nor the counterexample decides whether the same
$\exp\{c k\log(ed/k)\}$ lower bound holds for growing sparse $k$ in
either of these two regimes.

\begin{thebibliography}{9}
\bibitem{ABFR}
J. Aspnes, R. Beigel, M. Furst, and S. Rudich.
The expressive power of voting polynomials.
\emph{Combinatorica} 14(2):135--148, 1994.
\href{https://www.cs.yale.edu/homes/aspnes/papers/stoc91voting.pdf}{Author manuscript}.
\bibitem{flajolet2009}
P. Flajolet and R. Sedgewick. \emph{Analytic Combinatorics}.
Cambridge University Press, 2009. Chapter VI.
\bibitem{mai2021coresets}
T. Mai, C. Musco, and A. Rao.
Coresets for classification---simplified and strengthened.
\emph{Advances in Neural Information Processing Systems} 34, 2021.
\href{https://papers.neurips.cc/paper_files/paper/2021/file/6098ed616e715171f0dabad60a8e5197-Paper.pdf}{Primary paper}; Theorems 3 and 5.
\bibitem{tropp2008}
J. A. Tropp. Norms of random submatrices and sparse approximation.
\emph{Comptes Rendus Math\'ematique} 346:1271--1274, 2008.
\href{https://tropp.caltech.edu/papers/Tro08-Norms-Random-v3.pdf}{Author manuscript},
Theorem 1.1.
\bibitem{prekopa1973}
A. Pr\'ekopa. On logarithmic concave measures and functions.
\emph{Acta Scientiarum Mathematicarum} 34:335--343, 1973.
\end{thebibliography}
\end{document}
